% Options for packages loaded elsewhere
% Options for packages loaded elsewhere
\PassOptionsToPackage{unicode}{hyperref}
\PassOptionsToPackage{hyphens}{url}
\PassOptionsToPackage{dvipsnames,svgnames,x11names}{xcolor}
%
\documentclass[
  letterpaper,
  DIV=11,
  numbers=noendperiod]{scrartcl}
\usepackage{xcolor}
\usepackage{amsmath,amssymb}
\setcounter{secnumdepth}{5}
\usepackage{iftex}
\ifPDFTeX
  \usepackage[T1]{fontenc}
  \usepackage[utf8]{inputenc}
  \usepackage{textcomp} % provide euro and other symbols
\else % if luatex or xetex
  \usepackage{unicode-math} % this also loads fontspec
  \defaultfontfeatures{Scale=MatchLowercase}
  \defaultfontfeatures[\rmfamily]{Ligatures=TeX,Scale=1}
\fi
\usepackage{lmodern}
\ifPDFTeX\else
  % xetex/luatex font selection
\fi
% Use upquote if available, for straight quotes in verbatim environments
\IfFileExists{upquote.sty}{\usepackage{upquote}}{}
\IfFileExists{microtype.sty}{% use microtype if available
  \usepackage[]{microtype}
  \UseMicrotypeSet[protrusion]{basicmath} % disable protrusion for tt fonts
}{}
\makeatletter
\@ifundefined{KOMAClassName}{% if non-KOMA class
  \IfFileExists{parskip.sty}{%
    \usepackage{parskip}
  }{% else
    \setlength{\parindent}{0pt}
    \setlength{\parskip}{6pt plus 2pt minus 1pt}}
}{% if KOMA class
  \KOMAoptions{parskip=half}}
\makeatother
% Make \paragraph and \subparagraph free-standing
\makeatletter
\ifx\paragraph\undefined\else
  \let\oldparagraph\paragraph
  \renewcommand{\paragraph}{
    \@ifstar
      \xxxParagraphStar
      \xxxParagraphNoStar
  }
  \newcommand{\xxxParagraphStar}[1]{\oldparagraph*{#1}\mbox{}}
  \newcommand{\xxxParagraphNoStar}[1]{\oldparagraph{#1}\mbox{}}
\fi
\ifx\subparagraph\undefined\else
  \let\oldsubparagraph\subparagraph
  \renewcommand{\subparagraph}{
    \@ifstar
      \xxxSubParagraphStar
      \xxxSubParagraphNoStar
  }
  \newcommand{\xxxSubParagraphStar}[1]{\oldsubparagraph*{#1}\mbox{}}
  \newcommand{\xxxSubParagraphNoStar}[1]{\oldsubparagraph{#1}\mbox{}}
\fi
\makeatother


\usepackage{longtable,booktabs,array}
\usepackage{calc} % for calculating minipage widths
% Correct order of tables after \paragraph or \subparagraph
\usepackage{etoolbox}
\makeatletter
\patchcmd\longtable{\par}{\if@noskipsec\mbox{}\fi\par}{}{}
\makeatother
% Allow footnotes in longtable head/foot
\IfFileExists{footnotehyper.sty}{\usepackage{footnotehyper}}{\usepackage{footnote}}
\makesavenoteenv{longtable}
\usepackage{graphicx}
\makeatletter
\newsavebox\pandoc@box
\newcommand*\pandocbounded[1]{% scales image to fit in text height/width
  \sbox\pandoc@box{#1}%
  \Gscale@div\@tempa{\textheight}{\dimexpr\ht\pandoc@box+\dp\pandoc@box\relax}%
  \Gscale@div\@tempb{\linewidth}{\wd\pandoc@box}%
  \ifdim\@tempb\p@<\@tempa\p@\let\@tempa\@tempb\fi% select the smaller of both
  \ifdim\@tempa\p@<\p@\scalebox{\@tempa}{\usebox\pandoc@box}%
  \else\usebox{\pandoc@box}%
  \fi%
}
% Set default figure placement to htbp
\def\fps@figure{htbp}
\makeatother





\setlength{\emergencystretch}{3em} % prevent overfull lines

\providecommand{\tightlist}{%
  \setlength{\itemsep}{0pt}\setlength{\parskip}{0pt}}



 


\KOMAoption{captions}{tableheading}
\makeatletter
\@ifpackageloaded{caption}{}{\usepackage{caption}}
\AtBeginDocument{%
\ifdefined\contentsname
  \renewcommand*\contentsname{Table of contents}
\else
  \newcommand\contentsname{Table of contents}
\fi
\ifdefined\listfigurename
  \renewcommand*\listfigurename{List of Figures}
\else
  \newcommand\listfigurename{List of Figures}
\fi
\ifdefined\listtablename
  \renewcommand*\listtablename{List of Tables}
\else
  \newcommand\listtablename{List of Tables}
\fi
\ifdefined\figurename
  \renewcommand*\figurename{Figure}
\else
  \newcommand\figurename{Figure}
\fi
\ifdefined\tablename
  \renewcommand*\tablename{Table}
\else
  \newcommand\tablename{Table}
\fi
}
\@ifpackageloaded{float}{}{\usepackage{float}}
\floatstyle{ruled}
\@ifundefined{c@chapter}{\newfloat{codelisting}{h}{lop}}{\newfloat{codelisting}{h}{lop}[chapter]}
\floatname{codelisting}{Listing}
\newcommand*\listoflistings{\listof{codelisting}{List of Listings}}
\makeatother
\makeatletter
\makeatother
\makeatletter
\@ifpackageloaded{caption}{}{\usepackage{caption}}
\@ifpackageloaded{subcaption}{}{\usepackage{subcaption}}
\makeatother
\usepackage{bookmark}
\IfFileExists{xurl.sty}{\usepackage{xurl}}{} % add URL line breaks if available
\urlstyle{same}
\hypersetup{
  pdftitle={ISO/IEC 27001},
  pdfauthor={Joao Luiz D'Andrea},
  colorlinks=true,
  linkcolor={blue},
  filecolor={Maroon},
  citecolor={Blue},
  urlcolor={Blue},
  pdfcreator={LaTeX via pandoc}}


\title{ISO/IEC 27001}
\usepackage{etoolbox}
\makeatletter
\providecommand{\subtitle}[1]{% add subtitle to \maketitle
  \apptocmd{\@title}{\par {\large #1 \par}}{}{}
}
\makeatother
\subtitle{Information Security Management System --- Governance, Risk,
Controls, and Implementation Template}
\author{Joao Luiz D'Andrea}
\date{}
\begin{document}
\maketitle

\renewcommand*\contentsname{Table of contents}
{
\hypersetup{linkcolor=}
\setcounter{tocdepth}{3}
\tableofcontents
}

\section{Executive Summary}\label{executive-summary}

ISO/IEC 27001 provides the requirements for establishing, implementing,
maintaining, and continually improving an information security
management system (ISMS). It enables organisations to manage information
security risks systematically and to demonstrate that information
security is governed through a structured management system.

This document provides a practical template for defining the
organisation's ISO/IEC 27001 governance model, ISMS scope, risk
management approach, control structure, implementation roadmap, and
assurance model. It can be used to document how the organisation aligns
information security governance, risk treatment, controls, and continual
improvement into a coherent operating model.

\section{Overview}\label{overview}

\subsection{Purpose}\label{purpose}

This document provides a practical template for defining, implementing,
operating, and assuring an information security management system
aligned to ISO/IEC 27001.

\subsection{Objectives}\label{objectives}

The main objectives of applying ISO/IEC 27001 are to:

\begin{itemize}
\tightlist
\item
  establish a formal and auditable information security management
  system
\item
  protect confidentiality, integrity, and availability of information
\item
  align information security with business objectives and enterprise
  risk
\item
  apply systematic risk assessment and risk treatment
\item
  define and maintain appropriate information security controls
\item
  provide evidence of governance, monitoring, and continual improvement
\item
  support internal assurance and external certification where required
\end{itemize}

\subsection{Scope}\label{scope}

Describe the scope of ISO/IEC 27001 applicability in the organisation.

\subsubsection{In scope}\label{in-scope}

\begin{itemize}
\tightlist
\item
  the defined ISMS boundary
\item
  business units, processes, and services within scope
\item
  information assets and supporting technology
\item
  applications, infrastructure, platforms, and cloud services
\item
  people, facilities, and third parties relevant to the ISMS
\item
  security governance, risk treatment, and control monitoring activities
\item
  incident management, resilience, and continual improvement activities
  within the ISMS scope
\end{itemize}

\subsubsection{Out of scope}\label{out-of-scope}

Document justified exclusions, if any.

\subsection{Intended Audience}\label{intended-audience}

\begin{itemize}
\tightlist
\item
  Board and executive leadership
\item
  Chief Information Security Officer
\item
  Chief Information Officer
\item
  Enterprise Risk Management teams
\item
  Security governance, architecture, and operations teams
\item
  IT operations and service management teams
\item
  Business owners and control owners
\item
  Compliance and legal teams
\item
  Internal audit and assurance teams
\item
  PMO / transformation teams
\end{itemize}

\section{Standard Foundation}\label{standard-foundation}

\subsection{Background}\label{background}

ISO/IEC 27001 is the world's best-known standard for information
security management systems and specifies the requirements for
establishing, implementing, maintaining, and continually improving an
ISMS. It is designed to support a risk-based approach that can be scaled
to organisations of different sizes and sectors.
:contentReference{oaicite:1}

\subsection{ISMS Principles}\label{isms-principles}

Document how the organisation applies the core principles of the ISMS:

\begin{itemize}
\tightlist
\item
  risk-based thinking
\item
  leadership and accountability
\item
  information protection based on business needs
\item
  systematic control selection
\item
  evidence-based management
\item
  continual improvement
\item
  integration with organisational processes and governance
\end{itemize}

\subsection{Core Components}\label{core-components}

Describe the core components of the ISO/IEC 27001 model:

\begin{itemize}
\tightlist
\item
  organisational context
\item
  leadership and policy
\item
  planning
\item
  support
\item
  operation
\item
  performance evaluation
\item
  improvement
\item
  information security risk assessment
\item
  information security risk treatment
\item
  Statement of Applicability
\end{itemize}

\subsection{Relationship to ISO/IEC 27002 and Annex
A}\label{relationship-to-isoiec-27002-and-annex-a}

Describe how the organisation uses:

\begin{itemize}
\tightlist
\item
  ISO/IEC 27001 for management system requirements
\item
  Annex A as the information security controls reference
\item
  ISO/IEC 27002:2022 for guidance on information security controls
\end{itemize}

ISO states that ISO/IEC 27001 contains the ISMS requirements, while
ISO/IEC 27002 provides detailed control guidance. ISO material and SC 27
guidance also explain that Annex A is a normative reference set of
controls that must be considered as part of risk treatment and
comparison for completeness. :contentReference{oaicite:2}

\subsubsection{Example standards table}\label{example-standards-table}

\begin{longtable}[]{@{}
  >{\raggedright\arraybackslash}p{(\linewidth - 2\tabcolsep) * \real{0.5000}}
  >{\raggedright\arraybackslash}p{(\linewidth - 2\tabcolsep) * \real{0.5000}}@{}}
\toprule\noalign{}
\begin{minipage}[b]{\linewidth}\raggedright
Standard / Component
\end{minipage} & \begin{minipage}[b]{\linewidth}\raggedright
Purpose
\end{minipage} \\
\midrule\noalign{}
\endhead
\bottomrule\noalign{}
\endlastfoot
ISO/IEC 27001 & Defines ISMS requirements \\
Annex A & Reference set of information security controls \\
ISO/IEC 27002 & Guidance on the implementation and management of
controls \\
Statement of Applicability & Documents selected, excluded, and justified
controls \\
\end{longtable}

\section{Governance Model}\label{governance-model}

\subsection{Governance Objectives}\label{governance-objectives}

Define the governance objectives for ISO/IEC 27001 adoption:

\begin{itemize}
\tightlist
\item
  establish executive accountability for information security
\item
  define the ISMS scope and governance boundary
\item
  align security priorities with enterprise strategy and risk appetite
\item
  assign ownership for risk treatment and control performance
\item
  embed information security into operations, change, architecture, and
  supplier management
\item
  maintain evidence of implementation and operating effectiveness
\item
  support continual improvement through monitoring, review, corrective
  action, and audit
\end{itemize}

\subsection{Governance Structure}\label{governance-structure}

Describe the governance bodies and their responsibilities.

\subsubsection{Example governance
bodies}\label{example-governance-bodies}

\begin{itemize}
\tightlist
\item
  Board / Risk Committee
\item
  Executive Information Security Steering Committee
\item
  Information Security Forum
\item
  Enterprise Risk Committee
\item
  Architecture Review Board
\item
  Change Advisory Board
\item
  Internal Audit
\end{itemize}

\subsection{Three Lines Model}\label{three-lines-model}

Describe responsibilities across the three lines.

\subsubsection{First line}\label{first-line}

\begin{itemize}
\tightlist
\item
  ownership of business processes, assets, and operational controls
\item
  execution of security controls
\item
  risk treatment implementation
\item
  evidence production and remediation execution
\end{itemize}

\subsubsection{Second line}\label{second-line}

\begin{itemize}
\tightlist
\item
  information security governance
\item
  policy and standard setting
\item
  risk oversight
\item
  compliance monitoring
\item
  challenge and advisory support
\end{itemize}

\subsubsection{Third line}\label{third-line}

\begin{itemize}
\tightlist
\item
  internal audit
\item
  independent assurance
\item
  control effectiveness review
\end{itemize}

\section{Roles and Responsibilities}\label{roles-and-responsibilities}

\subsection{Key Roles}\label{key-roles}

Document the main ISO/IEC 27001-related roles.

\begin{itemize}
\tightlist
\item
  Board
\item
  CEO / Executive Management
\item
  Chief Information Security Officer
\item
  Chief Information Officer
\item
  Chief Risk Officer
\item
  Business Service Owner
\item
  Information Asset Owner
\item
  Security Governance Lead
\item
  Security Architect
\item
  Security Operations Lead
\item
  IT Operations Lead
\item
  Risk Manager
\item
  Supplier / Third-Party Risk Manager
\item
  Internal Audit
\end{itemize}

\subsection{RACI Matrix}\label{raci-matrix}

Use the following model and adapt it.

\begin{longtable}[]{@{}
  >{\raggedright\arraybackslash}p{(\linewidth - 18\tabcolsep) * \real{0.1000}}
  >{\raggedright\arraybackslash}p{(\linewidth - 18\tabcolsep) * \real{0.1000}}
  >{\raggedright\arraybackslash}p{(\linewidth - 18\tabcolsep) * \real{0.1000}}
  >{\raggedright\arraybackslash}p{(\linewidth - 18\tabcolsep) * \real{0.1000}}
  >{\raggedright\arraybackslash}p{(\linewidth - 18\tabcolsep) * \real{0.1000}}
  >{\raggedright\arraybackslash}p{(\linewidth - 18\tabcolsep) * \real{0.1000}}
  >{\raggedright\arraybackslash}p{(\linewidth - 18\tabcolsep) * \real{0.1000}}
  >{\raggedright\arraybackslash}p{(\linewidth - 18\tabcolsep) * \real{0.1000}}
  >{\raggedright\arraybackslash}p{(\linewidth - 18\tabcolsep) * \real{0.1000}}
  >{\raggedright\arraybackslash}p{(\linewidth - 18\tabcolsep) * \real{0.1000}}@{}}
\toprule\noalign{}
\begin{minipage}[b]{\linewidth}\raggedright
Activity
\end{minipage} & \begin{minipage}[b]{\linewidth}\raggedright
Board
\end{minipage} & \begin{minipage}[b]{\linewidth}\raggedright
Exec
\end{minipage} & \begin{minipage}[b]{\linewidth}\raggedright
CISO
\end{minipage} & \begin{minipage}[b]{\linewidth}\raggedright
CIO
\end{minipage} & \begin{minipage}[b]{\linewidth}\raggedright
Risk
\end{minipage} & \begin{minipage}[b]{\linewidth}\raggedright
Business Owner
\end{minipage} & \begin{minipage}[b]{\linewidth}\raggedright
Security Ops
\end{minipage} & \begin{minipage}[b]{\linewidth}\raggedright
Architecture
\end{minipage} & \begin{minipage}[b]{\linewidth}\raggedright
Audit
\end{minipage} \\
\midrule\noalign{}
\endhead
\bottomrule\noalign{}
\endlastfoot
Approve ISMS policy and scope & A & R & C & C & C & I & I & I & I \\
Maintain risk assessment methodology & I & A & R & C & C & I & I & I &
I \\
Approve Statement of Applicability & I & A & R & C & C & C & C & C &
I \\
Implement security controls & I & I & A & C & I & C & R & C & I \\
Conduct internal ISMS review & I & A & R & C & C & C & C & I & I \\
Manage security incidents & I & A & R & C & C & C & R & I & I \\
Track corrective actions & I & A & R & C & C & C & C & I & I \\
Perform independent assurance & I & I & I & I & I & I & I & I & R/A \\
\end{longtable}

\section{Policy Framework}\label{policy-framework}

\subsection{Core Policies}\label{core-policies}

List and describe the main ISMS-related policies.

\begin{itemize}
\tightlist
\item
  Information Security Policy
\item
  Risk Assessment and Risk Treatment Methodology
\item
  Access Control Policy
\item
  Information Classification and Handling Policy
\item
  Asset Management Standard
\item
  Cryptography Policy
\item
  Logging and Monitoring Standard
\item
  Vulnerability Management Standard
\item
  Incident Management Policy
\item
  Business Continuity / Disaster Recovery Policy
\item
  Supplier Security Policy
\item
  Secure Development Policy
\item
  Human Resources Security Policy
\end{itemize}

\subsection{Standards and Procedures}\label{standards-and-procedures}

Document the standards, procedures, control baselines, and operational
guidance supporting the ISMS.

\section{Core Processes}\label{core-processes}

\subsection{Context and Scope
Management}\label{context-and-scope-management}

Describe how the organisation determines and maintains:

\begin{itemize}
\tightlist
\item
  internal and external issues
\item
  interested parties and their requirements
\item
  the ISMS scope statement
\item
  interfaces and dependencies
\item
  applicability boundaries
\end{itemize}

\subsection{Risk Assessment}\label{risk-assessment}

Describe the process for:

\begin{itemize}
\tightlist
\item
  defining risk criteria
\item
  identifying information security risks
\item
  assessing likelihood and impact
\item
  evaluating risks against criteria
\item
  prioritising treatment actions
\item
  documenting results
\end{itemize}

\subsection{Risk Treatment}\label{risk-treatment}

Describe the process for:

\begin{itemize}
\tightlist
\item
  selecting treatment options
\item
  determining necessary controls
\item
  comparing selected controls with Annex A
\item
  documenting inclusions and exclusions
\item
  assigning owners and due dates
\item
  approving residual risk
\item
  tracking treatment implementation
\end{itemize}

\subsection{Statement of
Applicability}\label{statement-of-applicability}

Document how the organisation maintains the Statement of Applicability,
including:

\begin{itemize}
\tightlist
\item
  selected controls
\item
  excluded controls
\item
  reasons for inclusion or exclusion
\item
  implementation status
\item
  associated policies or procedures
\item
  accountable owners
\end{itemize}

\subsection{Operational Control
Management}\label{operational-control-management}

Describe the processes for:

\begin{itemize}
\tightlist
\item
  access management
\item
  change control
\item
  vulnerability management
\item
  logging and monitoring
\item
  supplier oversight
\item
  backup and recovery
\item
  incident handling
\item
  exception management
\end{itemize}

\subsection{Incident Management}\label{incident-management}

Describe the lifecycle for:

\begin{itemize}
\tightlist
\item
  event detection
\item
  triage
\item
  escalation
\item
  containment
\item
  eradication
\item
  recovery
\item
  evidence handling
\item
  lessons learned
\item
  corrective action tracking
\end{itemize}

\subsection{Performance Evaluation and
Improvement}\label{performance-evaluation-and-improvement}

Describe the process for:

\begin{itemize}
\tightlist
\item
  monitoring and measurement
\item
  internal audit
\item
  management review
\item
  nonconformity handling
\item
  corrective action
\item
  continual improvement
\end{itemize}

\section{Control Framework}\label{control-framework}

\subsection{Control Domains}\label{control-domains}

Organise the control environment using the Annex A themes.

\subsubsection{Example control domains}\label{example-control-domains}

\begin{itemize}
\tightlist
\item
  organisational controls
\item
  people controls
\item
  physical controls
\item
  technological controls
\item
  governance and oversight
\item
  risk management
\item
  operational security
\item
  supplier security
\item
  incident management
\item
  resilience and recovery
\item
  monitoring and assurance
\item
  continual improvement
\end{itemize}

\subsection{Sample Control Catalogue}\label{sample-control-catalogue}

\begin{longtable}[]{@{}
  >{\raggedright\arraybackslash}p{(\linewidth - 10\tabcolsep) * \real{0.1667}}
  >{\raggedright\arraybackslash}p{(\linewidth - 10\tabcolsep) * \real{0.1667}}
  >{\raggedright\arraybackslash}p{(\linewidth - 10\tabcolsep) * \real{0.1667}}
  >{\raggedright\arraybackslash}p{(\linewidth - 10\tabcolsep) * \real{0.1667}}
  >{\raggedright\arraybackslash}p{(\linewidth - 10\tabcolsep) * \real{0.1667}}
  >{\raggedright\arraybackslash}p{(\linewidth - 10\tabcolsep) * \real{0.1667}}@{}}
\toprule\noalign{}
\begin{minipage}[b]{\linewidth}\raggedright
Control ID
\end{minipage} & \begin{minipage}[b]{\linewidth}\raggedright
Control Name
\end{minipage} & \begin{minipage}[b]{\linewidth}\raggedright
Objective
\end{minipage} & \begin{minipage}[b]{\linewidth}\raggedright
Control Owner
\end{minipage} & \begin{minipage}[b]{\linewidth}\raggedright
Frequency
\end{minipage} & \begin{minipage}[b]{\linewidth}\raggedright
Evidence
\end{minipage} \\
\midrule\noalign{}
\endhead
\bottomrule\noalign{}
\endlastfoot
ISO-01 & ISMS policy approved & Ensure formal leadership direction &
CISO & Annual & Approved policy \\
ISO-02 & Risk assessment performed & Identify and evaluate security
risks & Risk Manager & Annual / change-driven & Risk assessment
records \\
ISO-03 & Statement of Applicability maintained & Document selected and
excluded controls & Security Governance Lead & Annual / change-driven &
SoA \\
ISO-04 & Internal audit completed & Validate ISMS conformance and
effectiveness & Audit Lead & Annual & Audit reports \\
ISO-05 & Management review performed & Ensure leadership oversight and
decisions & Executive Sponsor & Annual & Management review minutes \\
ISO-06 & Corrective actions tracked & Improve nonconformity resolution &
ISMS Manager & Ongoing & Action tracker \\
\end{longtable}

\section{Artefacts and Evidence}\label{artefacts-and-evidence}

\subsection{Required Artefacts}\label{required-artefacts}

Document the key ISO/IEC 27001 artefacts.

\begin{itemize}
\tightlist
\item
  ISMS scope statement
\item
  information security policy
\item
  risk assessment methodology
\item
  risk treatment methodology
\item
  risk register
\item
  risk treatment plan
\item
  Statement of Applicability
\item
  asset inventory
\item
  control standards and procedures
\item
  supplier security assessments
\item
  vulnerability and patch records
\item
  access review records
\item
  incident register
\item
  internal audit plan
\item
  internal audit reports
\item
  management review outputs
\item
  corrective action register
\item
  competency and awareness records
\item
  metrics and dashboard packs
\item
  certification audit materials where applicable
\end{itemize}

\subsection{Evidence Model}\label{evidence-model}

For each artefact define:

\begin{itemize}
\tightlist
\item
  owner
\item
  purpose
\item
  location
\item
  frequency of update
\item
  retention period
\item
  assurance mechanism
\end{itemize}

\section{Architecture and Technology
Implications}\label{architecture-and-technology-implications}

\subsection{Security by Design}\label{security-by-design}

Describe how ISMS requirements are embedded into:

\begin{itemize}
\tightlist
\item
  enterprise architecture
\item
  solution design
\item
  cloud architecture
\item
  network and infrastructure design
\item
  identity and access management
\item
  data protection architecture
\item
  software delivery pipelines
\item
  supplier integration patterns
\item
  operational monitoring design
\end{itemize}

\subsection{Technology Landscape View}\label{technology-landscape-view}

Document the security-relevant technology landscape:

\begin{itemize}
\tightlist
\item
  business services and critical processes
\item
  application and data estate
\item
  infrastructure and cloud services
\item
  identity and access ecosystem
\item
  security tooling landscape
\item
  logging and monitoring architecture
\item
  backup and recovery architecture
\item
  third-party technology dependencies
\end{itemize}

\subsection{Asset and Dependency
Mapping}\label{asset-and-dependency-mapping}

Include or reference:

\begin{itemize}
\tightlist
\item
  asset inventories
\item
  application dependency diagrams
\item
  network and trust zone diagrams
\item
  data flow diagrams
\item
  supplier dependency maps
\item
  resilience and recovery dependency maps
\end{itemize}

\section{Risk Management}\label{risk-management}

\subsection{Information Security Risk
Model}\label{information-security-risk-model}

Describe how information security risks are identified, assessed,
treated, accepted, escalated, and monitored within the ISMS.

\subsection{Example Risk Categories}\label{example-risk-categories}

\begin{itemize}
\tightlist
\item
  unclear ISMS scope or governance
\item
  incomplete asset visibility
\item
  weak identity and access management
\item
  poor control implementation
\item
  insecure change and development practices
\item
  inadequate monitoring and logging
\item
  supplier security exposure
\item
  insufficient incident response capability
\item
  weak recovery arrangements
\item
  poor evidence and assurance discipline
\end{itemize}

\subsection{Risk Register Template}\label{risk-register-template}

\begin{longtable}[]{@{}
  >{\raggedright\arraybackslash}p{(\linewidth - 16\tabcolsep) * \real{0.1111}}
  >{\raggedright\arraybackslash}p{(\linewidth - 16\tabcolsep) * \real{0.1111}}
  >{\raggedright\arraybackslash}p{(\linewidth - 16\tabcolsep) * \real{0.1111}}
  >{\raggedright\arraybackslash}p{(\linewidth - 16\tabcolsep) * \real{0.1111}}
  >{\raggedright\arraybackslash}p{(\linewidth - 16\tabcolsep) * \real{0.1111}}
  >{\raggedright\arraybackslash}p{(\linewidth - 16\tabcolsep) * \real{0.1111}}
  >{\raggedright\arraybackslash}p{(\linewidth - 16\tabcolsep) * \real{0.1111}}
  >{\raggedright\arraybackslash}p{(\linewidth - 16\tabcolsep) * \real{0.1111}}
  >{\raggedright\arraybackslash}p{(\linewidth - 16\tabcolsep) * \real{0.1111}}@{}}
\toprule\noalign{}
\begin{minipage}[b]{\linewidth}\raggedright
Risk ID
\end{minipage} & \begin{minipage}[b]{\linewidth}\raggedright
Risk Description
\end{minipage} & \begin{minipage}[b]{\linewidth}\raggedright
Cause
\end{minipage} & \begin{minipage}[b]{\linewidth}\raggedright
Impact
\end{minipage} & \begin{minipage}[b]{\linewidth}\raggedright
Likelihood
\end{minipage} & \begin{minipage}[b]{\linewidth}\raggedright
Inherent Risk
\end{minipage} & \begin{minipage}[b]{\linewidth}\raggedright
Control Strength
\end{minipage} & \begin{minipage}[b]{\linewidth}\raggedright
Residual Risk
\end{minipage} & \begin{minipage}[b]{\linewidth}\raggedright
Owner
\end{minipage} \\
\midrule\noalign{}
\endhead
\bottomrule\noalign{}
\endlastfoot
IS-01 & Example & Example & Example & Medium & High & Medium & Medium &
Example \\
\end{longtable}

\section{Metrics and Reporting}\label{metrics-and-reporting}

\subsection{Key Metrics}\label{key-metrics}

Examples of ISMS metrics:

\begin{itemize}
\tightlist
\item
  percentage of assets classified
\item
  percentage of risks assessed on schedule
\item
  percentage of treatment actions completed on time
\item
  percentage of controls implemented as planned
\item
  number of significant incidents
\item
  vulnerability remediation performance
\item
  access review completion rate
\item
  supplier assessment coverage
\item
  number of audit findings open vs closed
\item
  percentage of corrective actions closed on time
\item
  percentage of management review actions completed
\item
  percentage of ISMS objectives achieved
\end{itemize}

\subsection{Dashboard Structure}\label{dashboard-structure}

Suggested dashboard views:

\begin{itemize}
\tightlist
\item
  executive information security dashboard
\item
  ISMS performance dashboard
\item
  risk and treatment dashboard
\item
  incident and vulnerability dashboard
\item
  supplier security dashboard
\item
  audit and corrective action dashboard
\end{itemize}

\section{Audit and Assurance}\label{audit-and-assurance}

\subsection{Assurance Model}\label{assurance-model}

Document the assurance approach across:

\begin{itemize}
\tightlist
\item
  management self-assessment
\item
  control owner attestations
\item
  second-line compliance reviews
\item
  internal audit
\item
  certification audit support where applicable
\item
  corrective action verification
\end{itemize}

\subsection{Testing Themes}\label{testing-themes}

\begin{itemize}
\tightlist
\item
  scope and governance effectiveness
\item
  risk assessment and treatment quality
\item
  Statement of Applicability quality
\item
  control implementation and operating effectiveness
\item
  incident handling
\item
  supplier control coverage
\item
  management review effectiveness
\item
  corrective action discipline
\item
  evidence completeness and traceability
\end{itemize}

\section{Implementation Roadmap}\label{implementation-roadmap}

\subsection{Maturity Phases}\label{maturity-phases}

\subsubsection{Phase 1: Foundation}\label{phase-1-foundation}

\begin{itemize}
\tightlist
\item
  define organisational context and ISMS scope
\item
  establish governance and executive sponsorship
\item
  approve core policy framework
\item
  define risk assessment and treatment methods
\item
  build initial asset inventory and baseline controls
\end{itemize}

\subsubsection{Phase 2: Control Design}\label{phase-2-control-design}

\begin{itemize}
\tightlist
\item
  perform risk assessment
\item
  create risk treatment plan
\item
  develop Statement of Applicability
\item
  implement priority control improvements
\item
  establish incident, supplier, and monitoring processes
\end{itemize}

\subsubsection{Phase 3: Integration}\label{phase-3-integration}

\begin{itemize}
\tightlist
\item
  embed ISMS into change, architecture, procurement, and operations
\item
  improve evidence collection and monitoring
\item
  establish internal audit and management review cycle
\item
  align corrective action tracking with enterprise governance
\end{itemize}

\subsubsection{Phase 4: Optimisation}\label{phase-4-optimisation}

\begin{itemize}
\tightlist
\item
  refine metrics and dashboards
\item
  strengthen automation and control monitoring
\item
  improve management review quality
\item
  enhance assurance readiness and continual improvement
\item
  support certification maintenance where applicable
\end{itemize}

\subsection{Roadmap Table}\label{roadmap-table}

\begin{longtable}[]{@{}
  >{\raggedright\arraybackslash}p{(\linewidth - 12\tabcolsep) * \real{0.1429}}
  >{\raggedright\arraybackslash}p{(\linewidth - 12\tabcolsep) * \real{0.1429}}
  >{\raggedright\arraybackslash}p{(\linewidth - 12\tabcolsep) * \real{0.1429}}
  >{\raggedright\arraybackslash}p{(\linewidth - 12\tabcolsep) * \real{0.1429}}
  >{\raggedright\arraybackslash}p{(\linewidth - 12\tabcolsep) * \real{0.1429}}
  >{\raggedright\arraybackslash}p{(\linewidth - 12\tabcolsep) * \real{0.1429}}
  >{\raggedright\arraybackslash}p{(\linewidth - 12\tabcolsep) * \real{0.1429}}@{}}
\toprule\noalign{}
\begin{minipage}[b]{\linewidth}\raggedright
Initiative
\end{minipage} & \begin{minipage}[b]{\linewidth}\raggedright
Description
\end{minipage} & \begin{minipage}[b]{\linewidth}\raggedright
Owner
\end{minipage} & \begin{minipage}[b]{\linewidth}\raggedright
Priority
\end{minipage} & \begin{minipage}[b]{\linewidth}\raggedright
Timeline
\end{minipage} & \begin{minipage}[b]{\linewidth}\raggedright
Dependency
\end{minipage} & \begin{minipage}[b]{\linewidth}\raggedright
Status
\end{minipage} \\
\midrule\noalign{}
\endhead
\bottomrule\noalign{}
\endlastfoot
Example & Example & Example & High & Q1-Q2 & Example & Planned \\
\end{longtable}

\section{Cross-Framework Integration}\label{cross-framework-integration}

\subsection{Relationship to Other
Frameworks}\label{relationship-to-other-frameworks}

Describe how ISO/IEC 27001 integrates with:

\begin{itemize}
\tightlist
\item
  NIST Cybersecurity Framework
\item
  NIS2
\item
  DORA
\item
  GDPR
\item
  AI Act
\item
  COBIT
\item
  ITIL
\item
  TOGAF
\item
  enterprise risk management
\item
  business continuity and operational resilience frameworks
\end{itemize}

\subsection{Example Mapping Themes}\label{example-mapping-themes}

\begin{longtable}[]{@{}
  >{\raggedright\arraybackslash}p{(\linewidth - 4\tabcolsep) * \real{0.3333}}
  >{\raggedright\arraybackslash}p{(\linewidth - 4\tabcolsep) * \real{0.3333}}
  >{\raggedright\arraybackslash}p{(\linewidth - 4\tabcolsep) * \real{0.3333}}@{}}
\toprule\noalign{}
\begin{minipage}[b]{\linewidth}\raggedright
ISO/IEC 27001 Theme
\end{minipage} & \begin{minipage}[b]{\linewidth}\raggedright
Related Framework
\end{minipage} & \begin{minipage}[b]{\linewidth}\raggedright
Integration Point
\end{minipage} \\
\midrule\noalign{}
\endhead
\bottomrule\noalign{}
\endlastfoot
ISMS governance & COBIT & Governance objectives and accountability \\
Risk treatment & NIST CSF & Cyber risk management and capability
improvement \\
Incident management & NIS2 / DORA & Cyber incident governance and
reporting \\
Security architecture & TOGAF & Architecture governance and security
viewpoints \\
Privacy-related controls & GDPR & Data protection and processing
safeguards \\
\end{longtable}

\section{Maturity Model}\label{maturity-model}

\subsection{Suggested Maturity Levels}\label{suggested-maturity-levels}

\subsubsection{Level 1: Initial}\label{level-1-initial}

Security activities are reactive, fragmented, and poorly governed.

\subsubsection{Level 2: Repeatable}\label{level-2-repeatable}

Basic ISMS components exist but are inconsistently applied and weakly
evidenced.

\subsubsection{Level 3: Defined}\label{level-3-defined}

The ISMS scope, policies, risk treatment, and controls are documented
and broadly implemented.

\subsubsection{Level 4: Managed}\label{level-4-managed}

The ISMS is measured, reviewed, audited, and actively improved through
leadership oversight.

\subsubsection{Level 5: Optimised}\label{level-5-optimised}

Information security is fully integrated into enterprise governance,
architecture, operations, and supplier management with strong evidence
and continual improvement.

\subsection{Assessment Dimensions}\label{assessment-dimensions}

\begin{itemize}
\tightlist
\item
  governance
\item
  scope and context management
\item
  risk assessment
\item
  risk treatment
\item
  control implementation
\item
  incident management
\item
  supplier security
\item
  metrics and reporting
\item
  audit and management review
\item
  continual improvement
\end{itemize}

\section{Practical Templates}\label{practical-templates}

\subsection{Suggested Appendices}\label{suggested-appendices}

\begin{itemize}
\tightlist
\item
  Appendix A: ISMS scope statement template
\item
  Appendix B: Risk assessment template
\item
  Appendix C: Risk treatment plan template
\item
  Appendix D: Statement of Applicability template
\item
  Appendix E: Information security control catalogue
\item
  Appendix F: Internal audit checklist
\item
  Appendix G: Management review agenda template
\item
  Appendix H: ISO/IEC 27001 RACI matrix
\end{itemize}

\section{Customisation Notes}\label{customisation-notes}

This template should be tailored to reflect:

\begin{itemize}
\tightlist
\item
  organisational structure
\item
  business and regulatory context
\item
  ISMS scope boundary
\item
  critical services and information assets
\item
  sourcing and third-party model
\item
  cloud and technology estate
\item
  maturity level
\item
  assurance and certification expectations
\item
  operating model and budget constraints
\end{itemize}

\section{References}\label{references}

Add authoritative references here, for example:

\begin{itemize}
\tightlist
\item
  ISO/IEC 27001:2022
\item
  ISO/IEC 27001:2022/Amd 1:2024
\item
  ISO/IEC 27002:2022
\item
  ISO/IEC 27000 family guidance
\item
  internal information security policies and standards
\item
  enterprise risk framework
\item
  architecture and resilience standards
\end{itemize}




\end{document}
